\begin{table}[H]
    \centering
    \caption{Kinerja Koreksi Thorndike Case III pada Populasi Simulasi dengan Korelasi yang Menurun pada Skor Tinggi (SLODR)}
    \label{tab:slodr_case3}
    \begin{spacing}{1.0}
    \resizebox{\textwidth}{!}{%
    \begin{tabular}{ccccccccc}
        \toprule
        & \multicolumn{2}{c}{Korelasi Bersyarat} & \multicolumn{3}{c}{$\bar{R}$} & \multicolumn{3}{c}{$S_Z$} \\
        \cmidrule(lr){2-3} \cmidrule(lr){4-6} \cmidrule(lr){7-9}
        $b$ & $g=-1$ & $g=+1$ & Benar & Case III & Bias & Benar & Case III & Bias \\
        \midrule
        0,00 & 0,440 & 0,440 & 0,440 & 0,437 & -0,002 & 68,9 & 68,7 & -0,2 \\
        0,02 & 0,527 & 0,470 & 0,499 & 0,437 & -0,061 & 72,3 & 68,7 & -3,6 \\
        0,04 & 0,620 & 0,501 & 0,561 & 0,437 & -0,123 & 75,7 & 68,7 & -7,0 \\
        \bottomrule
    \end{tabular}%
    }
    \end{spacing}
    \vspace{2mm}
    \begin{flushleft}
    \footnotesize\textit{Catatan.} SLODR adalah \textit{Spearman's Law of Diminishing Returns}. Populasi Gaussian dengan $SD$ nasional tidak resmi dan muatan $m(g) = m_0 - b\,g$; sampel terseleksi lewat skor komposit (probabilitas seleksi logistik) dikalibrasi agar sama dengan sampel teramati (rerata skor komposit 657,8, $SD$ 53,8, rerata korelasi \textit{raw sample} 0,300). Karena sampel sama pada semua $b$, taksiran Case III identik; kolom Benar adalah nilai pada populasi simulasi. Bias = Case III $-$ Benar.
    \end{flushleft}
\end{table}
